\documentclass[12pt]{article}
\usepackage{amsmath}
\usepackage{enumitem}

\begin{document}

\begin{center}
\Large \textbf{SEA3004F -- Module 1 Examples}
\end{center}

\vspace{0.5cm}
\begin{enumerate}

\item Compute the mean pressure at the Earth’s surface given the total mass of the atmosphere 
$M_A = 5.26 \times 10^{18}\,\text{kg}$, gravity $g = 9.81\,\text{m s}^{-2}$ and the radius of the Earth 
$a = 6371\,\text{km}$.

\item Which of these equations is the right definition of specific humidity?

\begin{enumerate}[label=\alph*)]
\item $q = \dfrac{M_v}{M_d}$ with $M_v$ the mass of water vapour and $M_d$ the mass of dry air 
\item $q = \dfrac{\rho_v}{\rho_v + \rho_d}$
\item $q_s = \dfrac{e_s}{R_v T (\rho_v + \rho_d)}$ 
\end{enumerate}

\item What is the US Standard Atmosphere (1976)?

\item Calculate the density of water vapour over the desert with a temperature of $32^\circ\text{C}$ and a pressure of $4\,\text{mbar}$.

\item Assuming an exponential pressure and density dependence
\[
p = p_s e^{-z/H}
\]
with scale height $H = 7.5\,\text{km}$, estimate the heights in the atmosphere at which:

\begin{enumerate}[label=\alph*)]
\item the air density is equal to $1\,\text{kg m}^{-3}$
\item the pressure is equal to $1\,\text{hPa}$
\end{enumerate}

Use a surface pressure value $p_s = 1000\,\text{hPa}$ and a surface density $\rho_s = 1.25\,\text{kg m}^{-3}$.

\item Briefly explain Dalton’s law.

\item Give a brief description of the difference between Practical Salinity (PSS-78) and Absolute Salinity (TEOS-10).

\item Differentiate the net planetary heat budget and ocean heat budget.

\item How is ocean–air heat transfer affected by the presence of sea ice?

\item Which pressure level is used as reference for the computation of potential temperature in the ocean?

\item Define density anomaly and write down the formula for the linear equation of state for seawater.

\end{enumerate}

\end{document}